% Préambule du rapport : mise en page, polices, couleurs, environnements. Compilation : XeLaTeX (voir compiler.ps1).
\usepackage{fontspec}
\usepackage[french]{babel}
\usepackage{geometry}
\geometry{a4paper, top=24mm, bottom=24mm, left=25mm, right=22mm, headheight=14pt}

% Polices : Latin Modern, la police standard de LaTeX (même rendu que les rapports de TP du module).
\setmainfont{Latin Modern Roman}[SmallCapsFont={Latin Modern Roman Caps}]
\setsansfont{Latin Modern Sans}
\setmonofont{Latin Modern Mono}[Scale=1.02, BoldFont={lmmonolt10-bold.otf}]
\newcommand{\titres}{\rmfamily}

\usepackage{xcolor}
\definecolor{inpt}{HTML}{0B5A99}   % bleu des titres : bleu du logo INPT (007BC7) assombri pour la lecture
\definecolor{inptpale}{HTML}{EAF3FA}
\definecolor{encre}{HTML}{000000}
\definecolor{encre2}{HTML}{4A4A4A}
\definecolor{filet}{HTML}{000000}
\definecolor{papier}{HTML}{F5F7FA}
\definecolor{ok}{HTML}{1F6B3A}
\definecolor{ko}{HTML}{9B2C2C}
\definecolor{attente}{HTML}{8A5A00}
\color{encre}

\usepackage{microtype}
\usepackage{graphicx}
\graphicspath{{figures/}}
\usepackage{float}
\usepackage{booktabs}
\usepackage{longtable}
\usepackage{tabularx}
\usepackage{array}
\usepackage{rotating}
\usepackage{setspace}
\setstretch{1.08}
\usepackage{ragged2e}
\usepackage{amsmath}
\usepackage{etoolbox}
\usepackage{afterpage}

\usepackage{tikz}
\usetikzlibrary{positioning, arrows.meta, shapes.geometric, shapes.misc, shapes.multipart, fit, calc, backgrounds}
\usepackage{tcolorbox}
\usepackage{listings}

% Titres
\usepackage{titlesec}
\titleformat{\chapter}[hang]
  {\huge\bfseries\color{inpt}}
  {\thechapter.}
  {0.5em}
  {}
\titlespacing*{\chapter}{0pt}{-12pt}{16pt}
\titleformat{\section}{\Large\bfseries\color{inpt}}{\thesection}{0.8em}{}
\titleformat{\subsection}{\large\bfseries\color{inpt}}{\thesubsection}{0.8em}{}
\titleformat{\subsubsection}{\normalsize\bfseries\color{encre}}{}{0em}{}
\titlespacing*{\section}{0pt}{18pt}{8pt}
\titlespacing*{\subsection}{0pt}{14pt}{6pt}

% En-têtes et pieds de page
\usepackage{fancyhdr}
\pagestyle{fancy}
\fancyhf{}
\renewcommand{\headrulewidth}{0.4pt}
\renewcommand{\headrule}{{\color{filet}\hrule height\headrulewidth width\headwidth}}
\fancyhead[L]{\small\color{inpt}Projet SOA -- Gestion des stages}
\fancyhead[R]{\small\color{inpt}INPT -- \textsc{Module SOA}}
\fancyfoot[C]{\small\thepage}
\renewcommand{\chaptermark}[1]{\markboth{#1}{}}
\fancypagestyle{plain}{\fancyhf{}\renewcommand{\headrulewidth}{0pt}\renewcommand{\headrule}{}\fancyfoot[C]{\small\thepage}}

% Légendes (sans le paquet caption, absent de l'installation MiKTeX de base)
\makeatletter
\renewcommand{\fnum@figure}{\textsc{Figure}~\thefigure}
\renewcommand{\fnum@table}{\textsc{Tableau}~\thetable}
% « Figure 1 – texte », comme dans les rapports de TP
\long\def\@makecaption#1#2{\vskip\abovecaptionskip
  \sbox\@tempboxa{#1 -- #2}%
  \ifdim \wd\@tempboxa >\hsize #1 -- #2\par\else\global\@minipagefalse\hb@xt@\hsize{\hfil\box\@tempboxa\hfil}\fi
  \vskip\belowcaptionskip}
\makeatother

\usepackage[hidelinks, pdfusetitle]{hyperref}
\hypersetup{pdfauthor={Hiba Elamdaoui, Aya Zahri, Hibatollah Ait Ahmed}, pdftitle={Gestion intelligente des stages et des offres de recrutement}}
\urlstyle{tt}

% Code
\lstdefinestyle{code}{
  basicstyle=\ttfamily\footnotesize,
  columns=fullflexible, keepspaces=true,
  breaklines=true, breakatwhitespace=false, postbreak=\mbox{\textcolor{filet}{$\hookrightarrow$}\space},
  frame=leftline, framerule=1.6pt, rulecolor=\color{inpt}, framesep=8pt, xleftmargin=10pt,
  backgroundcolor=\color{papier},
  aboveskip=10pt, belowskip=10pt,
  showstringspaces=false, tabsize=2,
  commentstyle=\color{encre2}\itshape, stringstyle=\color{inpt}, keywordstyle=\color{encre}\bfseries,
}
\lstdefinelanguage{XMLfr}{morestring=[b]", morecomment=[s]{<!--}{-->}, alsoletter={:-}, sensitive=true}
\lstdefinelanguage{JSONfr}{morestring=[b]", sensitive=true}
\lstdefinelanguage{HTTPfr}{morestring=[b]", sensitive=true, morecomment=[l]{-->}}
\lstset{style=code}

% Encadrés
% (bibliothèques « skins » de tcolorbox non utilisées : elles exigent tikzfill, absent de MiKTeX de base)
\newenvironment{encadre}[1][]
  {\begin{tcolorbox}[colback=inptpale, colframe=inpt, boxrule=0pt, leftrule=2pt, arc=0pt, left=10pt, right=10pt, top=6pt, bottom=6pt]%
   \ifstrempty{#1}{}{{\titres\bfseries\small\color{inpt}#1}\par\smallskip}}
  {\end{tcolorbox}}

% Emplacement de capture d'écran : affiché tant que le fichier n'est pas dans figures/.
\newcommand{\emplacement}[2]{%
  \begin{tikzpicture}
    \node[draw=encre2, dashed, line width=0.8pt, fill=papier, minimum width=\linewidth-1pt, minimum height=#2, align=center, text width=0.8\linewidth]
      {{\titres\bfseries\color{ko}CAPTURE À FAIRE}\\[4pt]
       {\small\texttt{rapport/figures/\detokenize{#1}}}\\[2pt]
       {\footnotesize\color{encre2}Voir \texttt{docs/CAPTURES\_A\_FAIRE.md}}};
  \end{tikzpicture}}

% \capture{fichier}{hauteur de l'emplacement}{légende}{étiquette} : capture faite à la main (Bonita, SoapUI, Postman).
\newcommand{\capture}[4]{%
  \begin{figure}[htbp]\centering
    % L'emplacement vide reste compact (34 mm) ; la capture réelle prend la place nécessaire.
    \IfFileExists{figures/#1}{\includegraphics[width=\linewidth,height=0.45\textheight,keepaspectratio]{#1}}{\emplacement{#1}{34mm}}
    \caption{#3}\label{#4}
  \end{figure}}

% \ecran{fichier}{hauteur gardée en px}{hauteur totale en px}{largeur}{légende}{étiquette} : capture de l'interface (Playwright),
% coupée en haut de page quand elle est très haute.
\newcommand{\ecran}[6]{%
  \begin{figure}[H]\centering
    {\setlength{\fboxsep}{0pt}\color{filet}\fbox{\includegraphics[width=#4\linewidth, viewport=0 {\numexpr #3-#2\relax} 1360 #3, clip]{#1}}}
    \caption{#5}\label{#6}
  \end{figure}}

% Divers
% Lignes de tableau générées : \input n'est pas développable dans un tableau, la primitive \@@input l'est.
\makeatletter\newcommand{\lignes}[1]{\@@input #1 }\makeatother
\newcommand{\code}[1]{\texttt{#1}}
\newcommand{\regle}[1]{\textbf{\textcolor{inpt}{#1}}}
\newcommand{\statut}[1]{\textsf{\small #1}}
\setlength{\parskip}{4pt plus 1pt}
\setlength{\parindent}{0pt}
\setlength{\emergencystretch}{3em}  % évite les débordements dus aux termes en chasse fixe
\renewcommand{\arraystretch}{1.25}
\setlength{\LTpre}{8pt}
\setlength{\LTpost}{8pt}
% Listes plus serrées (sans enumitem)
\makeatletter
\def\@listi{\leftmargin\leftmargini \topsep 4pt \parsep 1pt \itemsep 2pt}
\let\@listI\@listi
\makeatother
